\chapter{周长、面积、体积}\label{ch:g6:measure}

篱笆有多长、田地有多大、
水箱能装多少水？这是三个不同的问题，
对应三个不同的量——\omterm{def:g6:measure:perimeter}{周长}、\omterm{def:g6:measure:area}{面积}、\omterm{def:g6:measure:volume}{体积}——
各有各的单位。把它们混淆是几何中最常见的错误；
本章将把它们彻底理清。

\section{长度与周长}

\begin{definition}[周长]\label{def:g6:measure:perimeter}
图形的\emph{周长}\index{周长}是其边界的总长度。
周长用长度单位度量：毫米（mm）、
厘米（cm）、米（m）、千米（km），换算关系如下：
\[
1 \text{ 千米} = 1000 \text{ 米}, \qquad
1 \text{ 米} = 100 \text{ 厘米}, \qquad
1 \text{ 厘米} = 10 \text{ 毫米}.
\]
\end{definition}

\begin{example}\label{ex:g6:measure:perimeter}
长为 $L$、宽为 $w$ 的长方形，其\omterm{def:g6:measure:perimeter}{周长}为
\[
P = L + w + L + w = 2 \times (L + w).
\]
对于长 $7$ 厘米、宽 $4$ 厘米的长方形：$P = 2 \times (7 + 4) = 2 \times 11 =
22$ 厘米。边长为 $c$ 的正方形，其\omterm{def:g6:measure:perimeter}{周长}为 $4c$。
\end{example}

\begin{proposition}[圆的周长]\label{prop:g6:measure:circle}
圆的\omterm{def:g6:measure:perimeter}{周长}（或称\emph{圆\omterm{def:g6:measure:perimeter}{周长}}）：\omterm{def:g3:shapes:shapes}{半径}为 $r$ 的圆，其\omterm{def:g6:measure:perimeter}{周长}为
\[
P = 2 \pi r ,
\]
其中 $\pi \approx 3.14$，对每个圆都是同一个数。
\end{proposition}
\admitted

\begin{example}
一个圆形池塘的\omterm{def:g3:shapes:shapes}{半径}为 $5$ 米。它的边界长
为 $2 \times \pi \times 5 = 10\pi \approx 31.4$ 米。尽量保留精确值
$10\pi$，只在最后才\omterm{def:g6:decimals:rounding}{四舍五入}。
\end{example}

\section{面积}

\begin{definition}[面积]\label{def:g6:measure:area}
图形的\emph{面积}\index{面积}度量的是它所覆盖的表面：
其内可容纳多少个单位正方形。单位有：厘米$^2$（边长为
$1$ 厘米的正方形）、米$^2$、千米$^2$，\dots\ 注意：
\[
1 \text{ 米}^2 = 100 \times 100 \text{ 厘米}^2 = 10\,000 \text{ 厘米}^2
\]
——一平方米是长为 $100$ 厘米、宽为 $100$ 厘米的正方形，因此
相邻两级单位的换算\omterm{def:g5:division:multiple}{倍数}是 $100$，而不是 $10$。
\end{definition}

\begin{omfigure}
\begin{tikzpicture}[scale=0.68]
  \draw[gray!50, thin] (0,0) grid (7,4);
  \draw[very thick, omDef] (0,0) rectangle (7,4);
  \fill[omThm!25] (0,3) rectangle (1,4);
  \draw[omThm, thick] (0,3) rectangle (1,4);
  \node[font=\small, omThm] at (2.6,3.5) {$1$ 个单位正方形};
  \node[font=\small] at (3.5,-0.5) {$7 \times 4 = 28$ 个单位正方形};
\end{tikzpicture}

{\small 长方形的\omterm{def:g6:measure:area}{面积}：$7$ 列，每列 $4$ 个单位正方形，
共 $7 \times 4 = 28$ 个正方形。这就是\omterm{def:g6:measure:area}{面积} $=$ 长
$\times$ 宽。}
\end{omfigure}

\begin{proposition}[基本面积公式]\label{prop:g6:measure:formulas}
\begin{center}
\renewcommand{\arraystretch}{1.5}
\begin{tabular}{l|l}
图形 & \omterm{def:g6:measure:area}{面积} \\
\hline
长方形（$L \times w$） & $A = L \times w$ \\
正方形（边长 $c$） & $A = c^2$ \\[4pt]
\omterm{def:g6:shapes:triangles}{直角三角形}（直角边 $a$、$b$） & $A = \dfrac{a \times b}{2}$ \\[4pt]
圆（\omterm{def:g3:shapes:shapes}{半径} $r$） & $A = \pi r^2$ \\
\end{tabular}
\end{center}
\end{proposition}

\begin{proof}[直角三角形的证明]
两个直角边分别为 $a$、$b$ 的\omterm{def:g6:shapes:triangles}{直角三角形}沿斜边
拼接成一个 $a \times b$ 的长方形。因此三角形的\omterm{def:g6:measure:area}{面积}是
长方形\omterm{def:g6:measure:area}{面积}的\omterm{def:g3:division:half}{一半}，即 $\frac{ab}{2}$。（长方形的公式即上文的单位正方形计数；
圆的公式在本阶段不加证明地承认，见
\cref{ch:g7:areas}。）
\end{proof}

\begin{omfigure}
\begin{tikzpicture}[scale=0.75]
  \draw[thick, fill=omDef!15] (0,0) -- (4,0) -- (0,2.5) -- cycle;
  \draw[thick, omExo, dashed] (4,0) -- (4,2.5) -- (0,2.5);
  \draw[thin] (0.3,0) -- (0.3,0.3) -- (0,0.3);
  \node[below, font=\small] at (2,-0.1) {$a$};
  \node[left, font=\small] at (0,1.25) {$b$};
  \node[font=\small, omExo] at (2.8,1.7) {第二份};
\end{tikzpicture}

{\small \omterm{def:g6:shapes:triangles}{直角三角形}是长方形的\omterm{def:g3:division:half}{一半}，因此公式为
$\frac{a\times b}{2}$。}
\end{omfigure}

\begin{example}[组合图形]\label{ex:g6:measure:composite}
一个 L 形房间是 $6$ 米 $\times$ $4$ 米的长方形，去掉一个 $2$ 米
$\times$ $2$ 米的正方形角。其\omterm{def:g6:measure:area}{面积}的计算步骤如下：
\begin{enumerate}
  \item 整个长方形：$6 \times 4 = 24$ 米$^2$；
  \item 去掉的正方形：$2 \times 2 = 4$ 米$^2$；
  \item 剩余\omterm{def:g6:measure:area}{面积}：$24 - 4 = 20$ 米$^2$。
\end{enumerate}
其\emph{\omterm{def:g6:measure:perimeter}{周长}}的数值 $20$ 不带\omterm{def:g6:measure:area}{面积}单位：沿 L 形绕行一周，
边界仍长 $6 + 4 + 6 + 4 = 20$ 米——切去一角的两道切口
只是替换了两段等长的墙。数值相同纯属巧合，
但一个是平方米，另一个是米！
\end{example}

\begin{remark}[周长相同，面积不同]
两个图形可以\omterm{def:g6:measure:perimeter}{周长}相同而\omterm{def:g6:measure:area}{面积}差别很大：一个
$5 \times 5$ 的正方形与 $9 \times 1$ 的长方形\omterm{def:g6:measure:perimeter}{周长}都是
$20$，而\omterm{def:g6:measure:area}{面积}分别为 $25$ 和 $9$。\omterm{def:g6:measure:perimeter}{周长}与\omterm{def:g6:measure:area}{面积}确实是
两个相互独立的量。
\end{remark}

\section{体积}

\begin{definition}[体积]\label{def:g6:measure:volume}
\omterm{def:g5:solids:def}{立体图形}的\emph{体积}\index{体积}度量它所占据的空间：
即其内部能放入多少个单位\omterm{def:g5:solids:def}{正方体}。单位有 cm$^3$、m$^3$、\dots\ 等，
且 $1$ m$^3 = 1\,000\,000$ cm$^3$（相邻两级相差
$1000$ 倍）。对于液体，还常用\emph{升}：
\[
1 \text{ L} = 1 \text{ dm}^3 = 1000 \text{ cm}^3,
\qquad
1 \text{ m}^3 = 1000 \text{ L}.
\]
\end{definition}

\begin{proposition}[长方体的体积]\label{prop:g6:measure:box}
一个\emph{\omterm{def:g5:solids:def}{长方体}}，长为 $L$、宽为 $w$、
高为 $h$，其\omterm{def:g6:measure:volume}{体积}为
\[
V = L \times w \times h .
\]
\end{proposition}

\begin{proof}
最底层包含 $L \times w$ 个单位\omterm{def:g5:solids:def}{正方体}
（见\cref{def:g6:measure:area}中的图，换成单位\omterm{def:g5:solids:def}{正方体}），
这样的层共有 $h$ 层。
\end{proof}

\begin{omfigure}
\begin{tikzpicture}[scale=0.6]
  \foreach \y in {0,1} {
    \foreach \x in {0,...,3} {
      \foreach \z in {0,1} {
        \begin{scope}[shift={({\x+0.45*\z},{\y+0.3*\z})}]
          \draw[fill=omDef!12] (0,0) rectangle (1,1);
          \draw[fill=omDef!20] (0,1) -- (0.45,1.3) -- (1.45,1.3) -- (1,1)
            -- cycle;
          \draw[fill=omDef!28] (1,0) -- (1.45,0.3) -- (1.45,1.3) -- (1,1)
            -- cycle;
        \end{scope}
      }
    }
  }
  \node[below, font=\small] at (2.5,-0.3)
    {$4 \times 2 \times 2 = 16$ 个单位正方体};
\end{tikzpicture}

{\small 装满单位\omterm{def:g5:solids:def}{正方体}的\omterm{def:g5:solids:def}{长方体}：每层 $4 \times 2$ 个\omterm{def:g5:solids:def}{正方体}，
共两层。}
\end{omfigure}

\begin{example}\label{ex:g6:measure:aquarium}
一个水族箱长 $60$ cm、宽 $30$ cm、高 $40$ cm。其\omterm{def:g6:measure:volume}{体积}为：
\[
V = 60 \times 30 \times 40 = 72\,000 \text{ cm}^3 = 72 \text{ dm}^3
= 72 \text{ L}.
\]
换算步骤如下：$1000$ cm$^3$ 为一升，
故 $72\,000 \div 1000 = 72$。
\end{example}

\section{练习}

\begin{exercise}[$\star$]\label{exo:g6:measure:1}
换算下列长度：$3.5$ m 化为 cm；$420$ mm 化为 cm；$0.8$ km 化为 m；
$25\,000$ m 化为 km。
\end{exercise}

\begin{exercise}[$\star$]\label{exo:g6:measure:2}
求下列图形的\omterm{def:g6:measure:perimeter}{周长}：长方形 $8$ cm $\times$ $3.5$ cm；
边长为 $6.2$ cm 的正方形；三边长分别为 $5$ cm、$7$ cm
和 $9$ cm 的三角形。
\end{exercise}

\begin{exercise}[$\star$]\label{exo:g6:measure:3}
一条圆形跑道的\omterm{def:g3:shapes:shapes}{半径}为 $50$ m。求一圈的长度
（先给出含 $\pi$ 的精确值，再\omterm{def:g6:decimals:rounding}{四舍五入}到米），需要跑多少圈
才能至少达到 $2$ km？
\end{exercise}

\begin{exercise}[$\star$]\label{exo:g6:measure:4}
求下列图形的\omterm{def:g6:measure:area}{面积}：长方形 $9$ cm $\times$ $4$ cm；边长为
$7$ m 的正方形；两直角边分别为 $6$ cm 和 $10$ cm 的\omterm{def:g6:shapes:triangles}{直角三角形}；\omterm{def:g3:shapes:shapes}{半径}为
$3$ cm 的圆面（先给出精确值，再\omterm{def:g6:decimals:rounding}{四舍五入}到 cm$^2$）。
\end{exercise}

\begin{exercise}[$\star$]\label{exo:g6:measure:5}
换算下列单位：$3$ m$^2$ 化为 cm$^2$；$45\,000$ cm$^2$ 化为 m$^2$；
$2.5$ L 化为 cm$^3$；$4\,500$ L 化为 m$^3$。
\end{exercise}

\begin{exercise}[$\star$]\label{exo:g6:measure:6}
一块长方形场地长 $120$ m、宽 $85$ m。围住它需要
多少米的篱笆？它的\omterm{def:g6:measure:area}{面积}是多少？
\end{exercise}

\begin{exercise}[$\star$]\label{exo:g6:measure:7}
求\omterm{def:g5:solids:def}{长方体} $5$ cm $\times$ $4$ cm $\times$ $10$ cm 的\omterm{def:g6:measure:volume}{体积}，
以及棱长为 $3$ cm 的\omterm{def:g5:solids:def}{正方体}的\omterm{def:g6:measure:volume}{体积}。
\end{exercise}

\begin{exercise}[$\star$]\label{exo:g6:measure:8}
画出两个\omterm{def:g6:measure:perimeter}{周长}为 $16$ cm 的不同长方形，并计算它们的
\omterm{def:g6:measure:area}{面积}。你所画的哪个长方形的\omterm{def:g6:measure:area}{面积}更大？
\end{exercise}

\begin{exercise}[$\star\star$]\label{exo:g6:measure:9}
一个 T 形图形由一个 $10 \times 2$ 的水平长方形
放在一个 $2 \times 6$ 的竖直长方形上方组成（单位：cm）。计算它的
\omterm{def:g6:measure:area}{面积}，再计算它的\omterm{def:g6:measure:perimeter}{周长}（沿边界仔细走一圈，
把每一条边都加上）。
\end{exercise}

\begin{exercise}[$\star\star$]\label{exo:g6:measure:10}
游泳池是一个\omterm{def:g5:solids:def}{长方体}，长 $25$ m，宽 $10$ m，
深 $2$ m。
\begin{enumerate}
  \item 池子装满水时，能装多少立方米的水？
  \item 这些水合多少升？
  \item 池子以每小时 $50\,000$ L 的速度注水，
        注满需要多长时间？
\end{enumerate}
\end{exercise}

\begin{exercise}[$\star\star$]\label{exo:g6:measure:11}
一个花园是边长为 $20$ m 的正方形，其中有一个圆形池塘，
池塘\omterm{def:g3:shapes:shapes}{半径}为 $3$ m。草地\omterm{def:g6:measure:area}{面积}是多少（先求精确值，再
\omterm{def:g6:decimals:rounding}{四舍五入}到整数 m$^2$）？
\end{exercise}

\begin{exercise}[$\star\star\star$]\label{exo:g6:measure:12}
一块巧克力条的尺寸为 $15$ cm $\times$ $6$ cm $\times$ $1$ cm。
制造商将\emph{三个}尺寸都扩大为原来的两倍。
\begin{enumerate}
  \item \omterm{def:g6:measure:volume}{体积}扩大为原来的多少倍？通过计算两个
        \omterm{def:g6:measure:volume}{体积}来验证。
  \item 价格变为原来的 $4$ 倍。大块巧克力是否
        比小块更划算？
\end{enumerate}
\end{exercise}

\section{问题：狄多女王的篱笆}

\begin{problem}[{周末问题——篱笆长度固定时，
哪种形状围出的土地最多？正方形胜过每一个长方形，
圆形胜过正方形，而谷仓的墙会
改变一切}]
\label{pb:g6:measure:1}
传说狄多女王在非洲海岸登陆后，
获准拥有“一张牛皮所能围住的土地”——于是她把牛皮
剪成一条极长极细的皮条，并围出足以
建立迦太基城的土地。现在这个问题交给你：一位
农夫恰好拥有 $20$ m 的篱笆。\Cref{exo:g6:measure:8}
已说明，\omterm{def:g6:measure:perimeter}{周长}相同的两个围栏可以有不同的
\omterm{def:g6:measure:area}{面积}；本题要找出\emph{最好的}围栏——并发现，
当谷仓的墙充当一条免费的边时，答案会
彻底改变。

\medskip
\noindent\textbf{第一部分——20 m 的篱笆。}
\begin{enumerate}
  \item 围栏必须是一个长方形，并恰好用完全部 $20$ m 篱笆。
        验证 $1 \times 9$ 与 $2 \times 8$ 的围栏都符合要求，
        并计算它们的\omterm{def:g6:measure:area}{面积}。
  \item 说明每个符合要求的围栏的长与宽
        之和都是 $10$ m。然后列出完整的表格，
        给出所有整数围栏（$1 \times 9$ 至 $5 \times 5$）
        及其\omterm{def:g6:measure:area}{面积}。哪一个最好？
  \item \omterm{def:g6:decimals:places}{小数}边长值得尝试吗？计算
        $4.5 \times 5.5$ 与 $4.9 \times 5.1$ 围栏的\omterm{def:g6:measure:area}{面积}，
        并与 $5 \times 5$ 的正方形比较。你
        有什么猜想？
  \item 第 2 题把篱笆问题化为一个纯数
        问题：\emph{在所有和为 $10$ 的数对中，
        哪一对的积最大？}根据你的表格作答，
        并陈述你观察到的一般规律。
  \item 下面用剪刀代替代数来说明，为什么离开正方形总会吃亏：
        从 $5 \times 5$ 的围栏出发，剪去一条窄条，
        再把另一条粘回去，变成 $6 \times 4$ 的围栏。
        哪一条被剪去，哪一条被粘上？
        这次交换为什么恰好损失一平方米？
        说明为什么再进一步（变为 $7 \times 3$）
        损失还要更多。
\end{enumerate}

\medskip
\noindent\textbf{第二部分——谷仓与圆。}
\begin{enumerate}[resume]
  \item 现在围栏靠着一堵长长的谷仓墙建造：这面墙
        代替围栏的一条边，$20$ m 的篱笆
        只围住其余三条边。列出整数围栏的表格
        （宽 $w$，沿墙的长 $L$，
        $w + L + w = 20$）及其\omterm{def:g6:measure:area}{面积}。现在哪一个围栏胜出——
        它是正方形吗？
  \item 用镜子来解释胜出的围栏
        （\cref{ch:g6:symmetry}）：以谷仓墙为\omterm{def:g6:symmetry:axis}{对称轴}作镜像，
        并考虑加倍后的围栏。加倍后的围栏用了多少篱笆？
        按第一部分的结论，哪一个加倍后的围栏最好？
        那么原来的围栏是什么样的？
  \item 狄多没有围成长方形。把 $20$ m 的篱笆
        弯成一个圆：取 $\pi \approx 3.14$，求它的
        \omterm{def:g3:shapes:shapes}{半径}（精确到 cm），再求它的\omterm{def:g6:measure:area}{面积}（精确到 m$^2$）
        （\cref{prop:g6:measure:circle}，
        \cref{prop:g6:measure:formulas}）。与
        $5 \times 5$ 的正方形比较。
  \item 反过来的问题：需要一个\omterm{def:g6:measure:area}{面积}恰好为 $36$ m$^2$ 的围栏，
        并且所用篱笆尽量少。比较
        长方形 $1 \times 36$、$2 \times 18$、$3 \times 12$、
        $4 \times 9$ 与 $6 \times 6$ 的\omterm{def:g6:measure:perimeter}{周长}：哪一种形状
        最省？这个问题与第一部分构成怎样的
        镜像关系？
  \item 用一句话说明：为什么罐子、管道和水箱常常做成
        \emph{圆形}？
\end{enumerate}

\medskip
\noindent\textbf{第三部分——\omterm{def:g6:measure:perimeter}{周长}与\omterm{def:g6:measure:area}{面积}互不相关。}
\begin{enumerate}[resume]
  \item 求一个\omterm{def:g6:measure:perimeter}{周长}超过 $100$ m、
        但\omterm{def:g6:measure:area}{面积}小于 $1$ m$^2$ 的长方形。（非常长、非常细。
        可以用\omterm{def:g6:decimals:places}{小数}。）
  \item 判断对错：“两个图形中，\omterm{def:g6:measure:perimeter}{周长}较大的
        那一个\omterm{def:g6:measure:area}{面积}也较大。”请给出
        本题中的一个反例。
  \item 取 $6 \times 4$ 的长方形，在一条长边的中间剪出一个 $1 \times 1$
        的正方形缺口（缺口向外开，
        像一颗缺失的牙齿）。计算有缺口图形的\omterm{def:g6:measure:area}{面积}
        与\omterm{def:g6:measure:perimeter}{周长}，并把两者都与原图形比较。
        两者各发生了什么变化？
  \item 制图者知道一个奇特的事实：地图越详细，
        海岸线测得的长度就\emph{越长}——每一次放大
        都会显露出新的小缺口，第 13 题说明了每个缺口
        的作用。用一两句话说明，为什么“布列塔尼海岸的长度”
        是一个难以捉摸的数，
        而“布列塔尼的\omterm{def:g6:measure:area}{面积}”却不是。
  \item 农夫的考题。有 $36$ m 的篱笆和可用的谷仓墙，
        求最好的长方形围栏（利用第 7 题的镜像），
        给出它的\omterm{def:g6:measure:area}{面积}，并与开阔地上用
        $36$ m 篱笆围成的最好围栏比较。谷仓的墙
        给农夫带来了多少好处？
\end{enumerate}
\end{problem}
